% !TEX root = ../math.tex
\chapter{Stimulus Optimization: Ideal Timer Placement}
\label{chap:optimize-stimulus-parameters}

\section{The Continuous One Bit Channel}
\subsection{Two Interpretations}

The continuous one bit channel can be conceived in two ways: as a continuous channel
with a one bit range, or as a series of times of identical events. The equivalency follows
from the fact that the continuous channel may be characterized by the times at which
its signal changes. To have a finite (differential) capacity we must constrain the channel to a minimum duration of signal,
or equivalently, a minimum time between events. For our purpose, we prefer the latter interpretation of
the one bit channel, since the user communicates using a sequence of identical events such as blinks
or taps.

\subsection{The Ideal Timing Distribution}
What then, is the capacity of the one bit channel? We prefer the event interpretation for analysis.
We consider the ideal distribution of the duration between events.
Let $T$ be the minimum duration between events and $X$ be the additional delay beyond the minimum,
so that the total duration between events is $T + X$.
Let $A = T + E[X]$ be the expected inter-arrival time.
The rate is given by
$$C = \max_{\mathcal{P}_X} \frac{h(X)}{A}.$$

We refer to the rate-maximizing distribution of $X$, $\mathcal{P}_X$, as the ideal timing distribution.
We will proceed to show that the ideal timing distribution is an exponential distribution.

We treat the capacity, expected value, and differential entropy as functionals of the probability density function $f(x)$.
To find the optimal distribution, we introduce a Lagrange multiplier $\mu$ for the normalization constraint $\int f(x) dx = 1$ and define the functional $J(f)$:
$$J(f) = \frac{h(X)}{A} - \mu \left( \int f(x)dx - 1 \right)$$
We maximize $J(f)$ by setting the variational derivative with respect to $f(x)$ to zero:
$$\frac{\delta J}{\delta f(x)} = 0$$
Using the quotient rule for variational derivatives, and noting that $\frac{\delta h(X)}{\delta f(x)} = -1 - \ln f(x)$ and $\frac{\delta A}{\delta f(x)} = \frac{\delta E[X]}{\delta f(x)} = x$, we have:
\[
\frac{\delta}{\delta f(x)} \left( \frac{h(X)}{A} \right) = \frac{A(-1 - \ln f(x)) - h(X)x}{A^2}
\]
Substituting this into the stationarity condition:
\[
\frac{A(-1 - \ln f(x)) - h(X)x}{A^2} - \mu = 0
\]
Solving for $f(x)$:
\begin{align*}
f(x) &\propto \exp\left( - \frac{h(X)}{A} x \right)
\end{align*}
proving that $X$ is exponentially distributed with rate parameter equal to the rate of the channel
(i.e., $\lambda^* = C$).
This fact lets us solve for the rate,
$$C = \frac{1 - \ln C}{T + 1/C}$$
which by algebra has the solution,
$$C = W(T) / T$$, where $W$ is the Lambert W function.
% TODO: Technically the argument to W should be dimensionless
% and we should write W(T/Δ)
\subsubsection{A Geometric Intuition}
The above result (that the ideal timing distribution is exponential) may
be appreciated more readily with a geometric argument.
We consider a large sample, $\mathbf{X}$, in $N$-dimensional space,
representing the sequence of additional delays $X_i$.
By the law of large numbers, its $L_1$ norm, $\|\mathbf{X}\|_1 = \sum_i X_i$, must be
tightly distributed about the $N$-simplex, $\|\mathbf{X}\|_1 = NE[X]$.
This can be thought of as the permissible region for $\mathbf{X}$. Thus the entropy maximizing
distribution for $\mathbf{X}$ should assign a uniform distribution over the $N$-simplex,
which is precisely what the exponential distribution does as it
depends only on the $L_1$ norm, $f(\mathbf{X}) \propto \exp(-\lambda \|\mathbf{X}\|_1)$.

(Incidentally, this conception also provides a geometric proof of the fact that the
differential entropy of the exponential distribution is $1 - \ln \lambda $, since 
the entropy of a uniform distribution over the $N$-simplex should be the logarithm of
its volume. Since its side length is $N/\lambda$, its volume is $(N/\lambda)^N / N!$,
and its entropy per sample (by Stirling's approximation) is
$$h(X) = \lim_{N \to \infty} \frac{1}{N} \left( N (\log(N - \log \lambda) - N (\log N - 1)) \right) = 1 - \ln \lambda$$
as expected.)




\section{Periodic Signals}
% Discuss maximizing the rate in practice
\subsection{Random Scan Times}
\subsection{Periodic signals}
% Discuss the advantages of making D a periodic variable
% - velocity is constant / doesn't need to be estimated (though this could be avoided)
% - info gain does not depend on the distribution of T
\subsection{Approximating A Uniform Distribution with a Gaussian Mixture}
UNFINISHED

We have considered the ideal timing distribution and the periodic timing distribution. 
How can we induce a uniform distribution on the user's signal?
There are several targets that the user may aim for, corresponding to all
the prefixes displayed.
Each target has an assigned signal time, $\mu_i$ and a fixed and known (mixing) probability, $\pi_i$. Given the noise $\mathcal{E} \sim \mathcal{N}(0, \sigma)$,
the distribution of the observed signal is a mixture of gaussians with means $mu_i$ and
constant variance $\sigma^2$.
Given the mixing probabilities, $\pi_i$, how should we position the $\mu_i$ to maximize the entropy of the signal?

Intuitively we want to space apart the $\mu_i$ to approximate a uniform distribution.
Those components with higher $pi_i$ should receive more space.

As soon as we receive a signal we recompute the posterior probabilities of the targets
and must assign new phases to them. Hence this computation is on the critical path
and ought to be computed in a few milliseconds. This rules out heavier techniques
like gradient descent on a numerical approximation of the entropy.

We discuss two heuristic approaches to the problem
as well as an approximation for the entropy which may computed efficiently.

\subsubsection{Two Heuristics}
UNFINISHED
We wish to space the components of the mixture as to approximate a uniform distribution.
One heuristic is to determine (by binary search) a value $C$ such that
$$\sum_i 2w_i = P$$
where the widths allotted to each component, $w_i$, are given by
$$w_i = \operatorname{Re}\left(\sigma \sqrt{2 \left( \ln \left( \frac{\pi_i}{\sqrt{2\pi}\sigma} \right) - C \right)}\right)$$
where $P$ is the length of the interval.
The means are then given by $\mu_i = w_i + \sum_{j=1}^{i-1} 2w_j$.

This heuristic is derived from approximating the entropy of the mixture as
$h(X) \approx - \ln f_{\min}$ where $f_{\min}$ is the minimum value of the mixture density,
and approximating the mixture by its greatest component at a given point,
$$f(x) \approx \max_i f_i(x)$$
The value of $C$ represents the minimum log density of the approximated mixture, i.e.,
the minimum over the interval of the maximum log density of the components.
The width $2w_i$ is the range over which the $i$th component may provide a log
density greater than $C$. If $\ln f_i(\mu_i) < C$, where $f_i(\mu_i)$ denotes the maximum density of the $i$th component at $x=0$, then clearly $w_i = 0$, otherwise it is as given above.
Our algorithm thus maximizes the minimum density of the maximum component over the interval.
Although somewhat crude, this heuristic is efficient.

A refinement may be made by instead approximating the mixture by the sum of its two greatest components,
and continuing to try to maximize the minimum density of the distribution.
% TODO: This is not technically accurate since a sub-C component could still be locally
% the greatest component.
In this case, adjacent components should be spaced apart such that the log of the sum of their densities
attains a minimum value of $C$.
Given two adjacent components, $i$ and $i+1$, we are interested in their separation, $l_i$.
Letting $\mu_i = 0$ and $\mu_{i+1} = l_i$, we may
solve for the minimum value of the two-component mixture:
$$f(x) = \pi_i \frac{1}{\sqrt{2 \pi \sigma^2}} \exp\left( - \frac{x^2}{2 \sigma^2} \right) + \pi_{i+1} \frac{1}{\sqrt{2 \pi \sigma^2}} \exp\left( - \frac{(x - l_i)^2}{2 \sigma^2} \right)$$
Note that completing the square for the sum of squared distances gives:
$$(x - \mu_1)^2 + (x - \mu_2)^2 = 2\left(x - \frac{\mu_1 + \mu_2}{2}\right)^2 + \frac{(\mu_1 - \mu_2)^2}{2}$$
Taking the derivative with respect to $x$ and setting it to zero we have,
$$\frac{d}{dx} f(x) = \frac{1}{\sqrt{2 \pi \sigma^2}} \left( \frac{-2x\pi_i}{2\sigma^2} \exp\left( - \frac{x^2}{2 \sigma^2} \right) - \frac{-2(x - l_i)\pi_{i+1}}{2\sigma^2} \exp\left( - \frac{(x - l_i)^2}{2 \sigma^2} \right) \right) = 0$$
Algebra yields,
$$\ln \left( \frac{R x}{l_i - x} \right) = \frac{x l_i}{\sigma^2} - \frac{l_i^2}{2 \sigma^2}$$,
where $R = \pi_i / \pi_{i+1}$.
Taken with the condition that $\ln f(x) = C$
the two variables, $C$ and $R$, implicitly determine $l_i$. So, using
a precomputed table of $l_i$ values for a range of $C$ and $R$ values, we can proceed as before,
using binary search to find the value of $C$ such that $\sum_i l_i = P$.
As before, we will exclude components with $\ln f_i(\mu_i) < C$ from calculations.
As long as each adjacent component has a maximum log density greater than $C$, it will be possible
to space them such that the log of their sum attains a minimum value of $C$.

% widths (max)
% widths (sum)
% cramer, renyi (quite simplex; approx ln z ~= z - 1)
% numerical integration (reasonable for low D)

\subsubsection{Approximations to the entropy based on pairwise distances}

In order to optimize the placement of the means to maximize entropy, we may minimize a proxy objective $J$ based on pairwise distances between the components. We consider approximations to the entropy for both Gaussian and von Mises mixture models. 

\paragraph{Gaussian Mixture Model with Rényi Entropy}
The Rényi 2-entropy, also known as collision entropy, is defined as $H_2(X) = -\ln \int p(x)^2 dx$. Maximizing this entropy is equivalent to minimizing the information potential $J = \int p(x)^2 dx$.
For a mixture of Gaussians $p(x) = \sum_{i=1}^N \pi_i \mathcal{N}(x | \mu_i, \sigma^2)$, the integral of the squared density is:
$$ J = \int \left( \sum_{i=1}^N \pi_i \mathcal{N}(x | \mu_i, \sigma^2) \right)^2 dx = \sum_{i=1}^N \sum_{j=1}^N \pi_i \pi_j \int \mathcal{N}(x | \mu_i, \sigma^2) \mathcal{N}(x | \mu_j, \sigma^2) dx $$
The integral of the product of two Gaussians is itself a Gaussian evaluated at the difference of their means:
$$ \int \mathcal{N}(x | \mu_i, \sigma^2) \mathcal{N}(x | \mu_j, \sigma^2) dx = \mathcal{N}(0 | \mu_i - \mu_j, 2\sigma^2) = \frac{1}{\sqrt{4\pi\sigma^2}} \exp\left( - \frac{(\mu_i - \mu_j)^2}{4\sigma^2} \right) $$
By dropping the leading scalar and the self-interaction terms ($i=j$) since they are constant with respect to the optimal placement of the means, the objective to minimize simplifies to the pairwise sum:
$$ J = \sum_{i < j} \pi_i \pi_j \exp\left( - \frac{(\mu_i - \mu_j)^2}{4\sigma^2} \right) $$

\paragraph{Gaussian Mixture Model with Kolchinsky Objective}
Kolchinsky and Tracey \cite{Kolchinsky_2017} introduced a family of computationally efficient estimators for mixture entropy that rely on calculating pairwise distances between component distributions, proving they form a strict lower bound on the true mixture entropy.
For homoskedastic Gaussian mixtures, maximizing this lower bound utilizing the Bhattacharyya distance yields a similar pairwise objective. The effective variance in the penalty term is simply doubled compared to the Rényi objective:
$$ J = \sum_{i < j} \pi_i \pi_j \exp\left( - \frac{(\mu_i - \mu_j)^2}{8\sigma^2} \right) $$

\paragraph{Von Mises Mixture Model with Rényi Entropy}
When the components are distributed on the unit circle as von Mises distributions with homoskedastic concentration $\kappa$, we again aim to minimize the integral of the squared density. The integral of the cross-term of two von Mises components evaluates to a modified Bessel function of the first kind of order zero, $I_0$. Stripping out constants and self-terms, the objective becomes:
$$ J = \sum_{i < j} \pi_i \pi_j I_0\left( 2\kappa \cos\left( \frac{\mu_i - \mu_j}{2} \right) \right) $$

\paragraph{Von Mises Mixture Model with Kolchinsky Objective}
Applying Kolchinsky's lower bound using the Bhattacharyya coefficient to the von Mises mixture model yields an objective with an identical structure to the Rényi case. The difference lies solely in the concentration parameter, which is halved:
$$ J = \sum_{i < j} \pi_i \pi_j I_0\left( \kappa \cos\left( \frac{\mu_i - \mu_j}{2} \right) \right) $$

\subsubsection{An approximation for the entropy via Cramér's theorem}
UNFINISHED
Here we discuss a closed form approximation for the entropy in terms of the
mixing weights, $\pi_i$, and the chosen means, $\mu_i$.
This enables efficient gradient descent on the means $\mu_i$.
We take a geometric approach and consider a large sample, $\mathbf{X}$,
in $N$-dimensional space.
Let $A_i$ be a discrete random variable taking on the values of $\mu_i$,
with respective probabilities $\pi_i$. Let $\mathcal{E}_i \sim \mathcal{N}(0, \sigma^2)$ be
random noise on the $i$th component. Then we have that 
$$\mathbf{X} = \mathbf{A} + \mathbf{\mathcal{E}}$$
And
\begin{align*}
    H(\mathbf{X}) &= H(\mathbf{A}, \mathbf{X}) - H(\mathbf{A} | \mathbf{X}) \\
    &= H(\mathbf{A}, \mathbf{\mathcal{E}}) - H(\mathbf{A} | \mathbf{X})  \quad \text{If $\mathbf{A}$ is fixed, $\mathbf{X}$ and $\mathbf{\mathcal{E}}$ differ only be a translation} \\
    &= H(\mathbf{\mathcal{E}}) + H(\mathbf{A}) - H(\mathbf{A} | \mathbf{X}) \\
    &= H(\mathbf{\mathcal{E}}) + I(\mathbf{A}; \mathbf{X})
\end{align*}
% NOTE: Since, A, is the transmitted signal
% TODO: we can think of this as positioning codeword centers
% (the spheres will differ in size because of the pi_i)
Since the entropy of the noise is constant, our objective is to maximize the mutual information
between $\mathbf{X}$ and the component choice, $\mathbf{A}$.
We frame this in terms of trying to infer the true component choice, $\mathbf{A}$, from the observed sample, $\mathbf{X}$.
To this end, we consider the prior distribution of the vertices' distances to the sample, $D = \|\mathbf{X} - \mathbf{B}\|_2^2$,
where $\mathbf{B}$ represents a candidate component choice and is distributed identically to $\mathbf{A}$.
% This needs to be A_i - B_i
Noting that, $D = \sum_i D_i$ where $D_i = (A_i + \mathcal{E}_i - B_i)^2$, we have that
$$P(\bar{D} = d) \approx \exp\left( - N I(d) \right)$$
and the likelihood for $\bar{D}$ is
$$P(\mathbf{X} | \bar{D} = d) = \exp\left( - \frac{1}{2\sigma^2} N d \right)$$
Thus the posterior mode of $\bar{D}$ is attained where
$N(I(d) + \frac{d}{2\sigma^2})$ is minimized,
which occurs when $I'(d) = - \frac{1}{2\sigma^2}$.
At this point, the prior probability of the posterior mode will be $\exp( - N I(d))$,
meaning that the observation delivers $N I(d)$ bits of information about the component choice,
or $I(d)$ bits per sample.

$I(d)$ had a closed form expression, although its derivation is a somewhat formidable piece of algebra.
We begin by defining the random variable $C_i = A_i - B_i$, which is also a discrete random variable
taking on the values of $d_{ij} = \mu_i - \mu_j$ with probabilities $\rho_{ij} = \pi_i \pi_j$.

Solving for $I$ in terms of its derivative, and setting it to $-\frac{1}{2\sigma^2}$, we have that

$$\mathbf{I\left(\theta = -\frac{1}{2\sigma^2}\right) = \frac{1}{2} \ln(2) - \frac{1}{4} - \ln(Z) - \frac{1}{8\sigma^2} \frac{Y}{Z}}$$

where

$$\mathbf{Z = \sum_{i,j} \rho_{ij} \exp\left( -\frac{d_{ij}^2}{4\sigma^2} \right)}$$

and

$$\mathbf{Y = \sum_{i,j} \rho_{ij} d_{ij}^2 \exp\left( -\frac{d_{ij}^2}{4\sigma^2} \right)}$$.

The final term may be recognized as the expected squared distance of $C_i$ under a Boltzmann distribution.


\section{User Parameters}
\subsection{The Likelihood Model}
% We identify two random variables in the user's behavior: $T$, the "scan time" required for
% the user to determine their next signal, and $X$ the duration between intervals itself.
We model our observations of the user's signal as being subject to two sources of noise: an outlier probability, $\rho$,
that the signal was unintended, and additive noise on intended signals.
We model the unintended signals as uniformly distributed, and the noise as a Gaussian. (We do not require that the noise have
zero mean, since the user may be consistently late or early, relative to the intended time.) Hence our likelihood model for the
received signal is a uniform-gaussian mixture model, characterized by the outlier probability, 
$\rho$, and the parameters of the noise, $\mu_\epsilon$, $\sigma_\epsilon$.

\subsection{Directional Statistics}
% Should be a uniform-von_mises mixture model
% which we can estimate using the EM algorithm
% which consists in soft assignment / assignment of responsibilities


\subsection{The Cost of Unintended Signals}
UNFINISHED
% TODO: properly write this as I(D; h)


The user's speed can be calculated from their performance parameters.
Here we consider the effect of $\rho$, the probability of
unintended signals.
Given an observation, $D$, the gain in logits for the true hypothesis
$h$ is $\ln P(D | h) - \ln P(D | \neg h)$. When the posterior probability of
$h$ is small, this gain is equal to the gain in the log posterior probability
of $h$.
Let $r$ denote the hypothesis that the signal was unintended, so that
$P(r) = \rho$. 

We state the following facts:
\begin{align*}
P(D | \neg h, r) = P(D | \neg h, \neg r) &= P(D | \neg h) \\
P(D | \neg h) &\approx P(D), \quad \text{when } P(h) \approx 0 \\
P(D | h, r) &= P(D),  \quad \text{assuming that $D$, $h$ are independent conditioned on $r$} \\
\end{align*}

The likelihood under the true hypothesis is
$$P(D | h) = \rho P(D | h, r) + (1 - \rho) P(D | h, \neg r)$$

We consider the case where $P > > \sigma$
so that the components of the mixture are well-separated.
This means that the contribution of the false component is negligible
and that we can approximate the likelihood $P(D|h)$ piecewise as follows:

\[
P(D|h) \approx 
\begin{cases}
P(D | h, r) P(r) & \text{if } D \sim P(D|h, r) \\
P(D | h, \neg r)P(\neg r) & \text{if } D \sim P(D|h, \neg r).
\end{cases}
\]

That is, when $D$ is likely drawn from the unintended distribution, $P(D|h) \approx P(D|h, r) P(r)$. When $D$ is likely drawn from the intended distribution, $P(D|h) \approx P(D|h, \neg r)P(\neg r)$.




\begin{align*}
E_{D | r, h} \left[ \ln P(D | h) \right] &\approx E_{D | r, h} \left[ \ln \Big( P(D | h, r) P(r) \Big) \right] \\
&= \ln \rho - h(D | h, r) \\
\end{align*}
and similarly
\begin{align*}
E_{D | \neg r, h} \left[ \ln P(D | h) \right] &\approx \ln (1 - \rho) - h(D | h, \neg r)
\end{align*}
The expected log likelihood of the data under the true hypothesis is
\begin{align*}
E_D \left[ \ln P(D | h) \right] &= \rho E_{D | r, h} \left[ \ln P(D | h) \right] + (1 - \rho) E_{D | \neg r, h} \left[ \ln P(D | h) \right] \\
&\approx \rho (\ln \rho - h(D | h, r)) + (1 - \rho) (\ln (1 - \rho) - h(D | h, \neg r)) \\
&= - H_2(\rho) - \rho h(D | h, r) - (1 - \rho) h(D | h, \neg r) \\
&= - H_2(\rho) - \rho h(D) - (1 - \rho) h(D | h, \neg r), \quad \text{assuming $D|r$, $h|r$, independent}
\end{align*}
Where $H_2$ is the binary entropy function.

The likelihood of the alternative hypothesis is
% separating cases here enlightens nothing
\begin{align*}
E_{D | h}\left[ \ln P(D | \neg h) \right] &\approx E_{D | h} \left[ \ln P(D) \right] \\
\end{align*}
To simplify this term we use its expectation over $h$.
Recalling the fact that,
$$ E_a \left[ E_{b | a} [f(b)] \right] = E_b [f(b)] $$,
we find that this term's expectation over $h$ is $-h(D)$.

So the total expected gain in the logits of our hypothesis is
\begin{align*}
E_{D | h}[\ln P(D | h) - P(D | \neg h)] &\approx - H_2(\rho) + (1 - \rho) h(D) - (1 - \rho) h(D | h, \neg r) \\
\end{align*}
In our particular model, $D$ is uniformly distributed over an interval of length $P$ and
the entropy of $D$ conditional on $h$ is just the entropy of the gaussian noise.
Denoting the latter term by $H_{\sigma}$, we have,

\begin{align*}
E_{D | h}[\ln P(D | h) - P(D | \neg h)] &\approx - H_2(\rho) + (1 - \rho) (\ln P - H_{\sigma})
\end{align*}

This has a nice interpretation: we gain $\ln P - H_{\sigma}$ bits of information
per intended signal, but the presence of unintended signals makes us pay a penalty in two ways:
only $(1 - \rho)$ of the time is our signal intended, and additionally we must
expend $H_2(\rho)$ bits to learn whether the signal was intended.

Unfortunately the singularity of the binary entropy function
at $\rho = 0$ becomes a practical problem for fast text entry. For example, a
seemingly low outlier rate of $\rho = 0.03$ imposes the relatively
high cost of $H_2(0.03) \approx 0.19$ bits in addition to the
$3\%$ loss that would be naively expected.
Under reasonable parameters for an intermediate user (e.g. $\sigma = 25ms$, $P = 700ms$),
these combined terms amount to a loss of more than $10\%$ of the channel's capacity.
Thus to achieve expert entry speeds, it is necessary to
make unintended signals very rarely.

